% 中文关键字
\newcommand{\keywordsCn}[1][图像融合；多曝光融合；视觉语言模型；深度学习；跨注意力机制；Restormer]{#1}
% 中文摘要
\begin{abstract}
\noindent\hspace*{2\ccwd}多曝光图像融合技术旨在将同一场景在不同曝光条件下拍摄的多幅图像合成为一幅包含完整场景细节的高质量图像。由于相机传感器动态范围的限制，单次曝光往往无法同时捕获场景中高亮区域和暗部区域的细节信息。通过融合不同曝光级别的图像，可以在标准显示设备上呈现更接近人眼真实视觉感知的图像效果，在数字摄影、安防监控、自动驾驶等领域具有广泛的应用价值。

近年来，基于深度学习的图像融合方法取得了显著的性能提升。然而，现有的深度学习方法大多仅依赖于源图像的视觉特征进行融合决策，缺乏对图像语义内容的深层理解。多模态研究的兴起为图像融合领域提供了新的研究思路。大语言模型在理解图像语义内容方面展现出了强大的能力，将文本语义信息引入图像融合过程，可以为融合决策提供高层语义指导。

针对上述问题，本文以视觉语言图像融合方法FILM（Image Fusion via Vision-Language Model）为基础，面向多曝光图像融合任务开展模型复现、推理流程适配、实验验证与交互式系统实现。本文的主要研究内容和成果如下：

（1）复现并适配了基于视觉语言模型的图像融合方法FILM。该方法利用预训练视觉语言模型生成的文本语义特征，通过跨注意力机制将文本特征与视觉特征进行对齐和融合，实现文本语义对视觉特征处理的引导。

（2）分析并实现了基于Restormer Transformer的跨注意力融合块（restormer\_cablock）。该模块对两幅源图像分别进行独立的特征提取和文本投影，以文本特征作为查询、图像特征作为键和值执行多头跨注意力计算，通过注意力池化和L1归一化实现文本语义引导下的视觉特征选择和重构。模型采用三级级联结构，逐级深化文本语义引导下的特征融合。

（3）在MEFB数据集上进行了实验验证。定量实验结果表明，FILM方法在标准差、空间频率、平均梯度、视觉信息保真度和边缘信息保持度五项客观评价指标上优于传统融合方法，信息熵指标与均值融合方法持平但低于部分传统方法。实验结果说明视觉语言语义特征能够在多曝光图像融合中提升细节保持和视觉信息保真能力。

（4）设计并实现了一个基于Web的交互式多曝光图像融合系统。系统采用前后端分离的B/S架构，前端基于Vue 3框架，后端基于FastAPI框架，支持AI融合和传统融合算法，提供图像上传、算法选择、结果对比查看、质量评估和历史记录管理等交互功能。系统测试结果表明各项功能运行正常，性能满足实际应用需求。

本文的研究工作以已有视觉语言图像融合方法为基础，将其适配到多曝光图像融合系统中，为深度学习图像融合算法的工程化应用提供了实践参考。
\end{abstract}

% 英文关键字
\newcommand{\keywordsEn}[1][Image Fusion; Multi-Exposure Fusion; Vision-Language Model; Deep Learning; Cross-Attention Mechanism; Restormer]{#1}
% 英文摘要
\begin{abstractEn}
\noindent\hspace*{1em}Multi-Exposure Image Fusion (MEF) technology aims to synthesize multiple images of the same scene captured under different exposure conditions into a single high-quality image containing complete scene details. Due to the limited dynamic range of camera sensors, a single exposure often fails to simultaneously capture details in both highlight and shadow regions of a scene. By fusing images of different exposure levels, it is possible to present image effects closer to the real visual perception of the human eye on standard display devices, which has broad application value in digital photography, security surveillance, autonomous driving, and other fields.

In recent years, deep learning-based image fusion methods have achieved significant performance improvements. However, most existing deep learning approaches rely solely on visual features from source images for fusion decisions, lacking deep understanding of the semantic content of images. The emergence of Vision-Language Models (VLMs) provides a new research direction for the image fusion field. Large Language Models have demonstrated powerful capabilities in understanding the semantic content of images. Introducing text semantic information into the image fusion process can provide high-level semantic guidance for fusion decisions.

To address the above issues, this thesis takes FILM (Image Fusion via Vision-Language Model) as the foundation and conducts model reproduction, inference adaptation, experimental validation, and interactive system implementation for the multi-exposure image fusion task. The main research contents and contributions of this thesis are as follows:

(1) The vision-language model-based image fusion method FILM is reproduced and adapted. This method utilizes text semantic features generated by a pre-trained vision-language model and aligns and fuses text features with visual features through a cross-attention mechanism, achieving text semantic guidance of visual feature processing.

(2) A Restormer Transformer-based cross-attention fusion block (restormer\_cablock) is analyzed and implemented. This module performs independent feature extraction and text projection for two source images respectively, executes multi-head cross-attention computation with text features as queries and image features as keys and values, and achieves text semantic-guided visual feature selection and reconstruction through attention pooling and L1 normalization. A three-stage cascaded structure is adopted to progressively deepen text semantic-guided feature fusion.

(3) Experimental validation is conducted on the MEFB dataset. Quantitative results show that the FILM method outperforms traditional fusion methods on five objective metrics, including standard deviation, spatial frequency, average gradient, visual information fidelity, and edge information preservation. The entropy metric is comparable to average fusion but lower than some traditional methods. The results indicate that vision-language semantic features can improve detail preservation and visual information fidelity in multi-exposure image fusion.

(4) A web-based interactive multi-exposure image fusion system is designed and implemented. The system adopts a B/S architecture with frontend-backend separation. The frontend is built on the Vue 3 framework, and the backend is based on the FastAPI framework. It supports AI fusion and traditional fusion algorithms, providing interactive functions including image upload, algorithm selection, result comparison, quality evaluation, and history record management. System testing results show that all functions operate normally and performance meets practical application requirements.

Based on an existing vision-language image fusion method, this thesis adapts it to a multi-exposure image fusion system and provides practical reference for the engineering application of deep learning-based image fusion algorithms.
\end{abstractEn}
